% Part: turing-machines
% Chapter: undecidability
% Section: universal-tm

\documentclass[../../../include/open-logic-section]{subfiles}

\begin{document}

\olfileid{tur}{und}{uni}
\olsection{సార్వత్రిక ట్యూరింగ్ యంత్రాలు}

\olref[enu]{sec}లో ప్రతి ట్యూరింగ్ యంత్రాన్నీ పూర్ణ సంఖ్యల పరిమిత క్రమంతో
ఎలా వర్ణించవచ్చో చర్చించాం. ఈ క్రమం ట్యూరింగ్ యంత్రపు స్థితులు, వర్ణమాల,
ప్రారంభ స్థితి, నిర్దేశాలను సంకేతీకరిస్తుంది. ఈ వర్ణనలన్నిటి సమితి
\tetoken{లెక్కించదగిన}{!!{enumerable}} అని కూడా పేర్కొన్నాం. అలాంటి వర్ణనల సమితి \tetoken{అనంతంగా లెక్కించదగిన}{!!{denumerable}}
కనుక, $\Nat$ నుండి ఈ వర్ణనలకు వెళ్లే \tetoken{ఒక సంగ్రస్త}{!!a{surjective}} ప్రమేయం ఉందని దీని
అర్థం. ఉదాహరణకు, కాంటర్ జిగ్‌జాగ్ పద్ధతితో అలాంటి \tetoken{ఒక సంగ్రస్త}{!!a{surjective}} ప్రమేయాన్ని
పొందవచ్చు. ట్యూరింగ్ యంత్రాల (వర్ణనల)న్నిటినీ లెక్కించడానికి అది ఒక పద్ధతిని
ఇస్తుంది. అలాంటి ఒక లెక్కింపును స్థిరపరిస్తే, $1$వ, $2$వ, \dots, $e$వ
ట్యూరింగ్ యంత్రం గురించి మాట్లాడడం అర్థవంతమవుతుంది. ఈ సంఖ్యలను
\emph{సూచికలు} అంటారు.

\begin{defn}
$M$ మన స్థిర లెక్కింపులో $e$వ ట్యూరింగ్ యంత్రం అయితే, $e$ను $M$ యొక్క
\emph{సూచిక} అంటాం. $e$వ ట్యూరింగ్ యంత్రాన్ని $M_e$ అని వ్రాస్తాం.
\end{defn}

ఒక యంత్రానికి ఒకటికంటే ఎక్కువ సూచికలు ఉండవచ్చు. ఉదాహరణకు, $M$ నిర్దేశాలను
జాబితా చేసే క్రమంలో దాని రెండు వర్ణనలు వేరుగా ఉండవచ్చు; ఆ వేర్వేరు వర్ణనలకు
వేర్వేరు సూచికలు ఉంటాయి.

ముఖ్యంగా, ఒక సూచిక నుండి దాని యంత్రం~$M$ యొక్క వర్ణనను ప్రభావకంగా గణించగలిగేలా,
అలాగే యంత్రం~$M$ వర్ణన నుండి దాని ఒక సూచికను ప్రభావకంగా గణించగలిగేలా,
ట్యూరింగ్ యంత్ర వర్ణనల లెక్కింపును ఇవ్వవచ్చు. అప్పుడు చర్చ్--ట్యూరింగ్
సిద్ధాంతప్రతిపాదన ప్రకారం, సూచిక~$e$ గల ట్యూరింగ్ యంత్రపు వర్ణనను తిరిగి
పొంది, అనుగుణమైన వర్ణనను అవుట్‌పుట్‌గా తన టేపుపై వ్రాసే ట్యూరింగ్ యంత్రాన్ని
కనుగొనవచ్చు. ఆ వర్ణన~$\TMstroke$ల దిమ్మల క్రమంగా ఉంటుంది (అవి $M_e$ను
వర్ణించే క్రమంలోని ధన పూర్ణ సంఖ్యలకు ప్రాతినిధ్యం వహిస్తాయి).

ఇది ఇచ్చినప్పుడు, కింది ప్రశ్నలు సహజమవుతాయి: ట్యూరింగ్ యంత్ర సూచికల ఏ
ప్రమేయాలు ట్యూరింగ్ యంత్రాలతో గణనీయాలు? ట్యూరింగ్ యంత్ర సూచికల ఏ లక్షణాలను
ట్యూరింగ్ యంత్రాలతో నిర్ణయించవచ్చు? ఒక ఉదాహరణ: సూచిక~$e$ను, సూచిక~$e$ గల
ట్యూరింగ్ యంత్రపు స్థితుల సంఖ్యకు పంపే ప్రమేయాన్ని ట్యూరింగ్ యంత్రం గణించగలదు.
అలాంటి ట్యూరింగ్ యంత్రం చేసే పని ఇదీ: $e$~$\TMstroke$ల ఒకే దిమ్మ గల టేపుపై
ప్రారంభమై, మొదట $e$ను దాని వర్ణనగా విసంకేతీకరిస్తుంది. అప్పుడు ఆ వర్ణన
టేపుపై~$\TMstroke$ల దిమ్మల క్రమంతో సూచించబడుతుంది. ఈ క్రమంలోని మొదటి
\tetoken{మూలకం}{!!{element}} స్థితుల సంఖ్య. కాబట్టి ఇప్పుడు చేయవలసిందల్లా మొదటి
$\TMstroke$ల దిమ్మ తప్ప మిగతాదంతా చెరిపి, తరువాత ఆగడం.\sourcecorrection{OLTETURUND-004}{ఆంగ్ల మూలం ``Since the first element in this sequence is the number of states. So~\dots'' అని కారణవాచకంతో మొదలైన వాక్యాన్ని అసంపూర్ణంగా వదిలింది. మొదటి మూలకమే స్థితుల సంఖ్య అనే పూర్తి వాక్యంగా ఇచ్చి, దాని నుండి తరువాతి చర్యను అనుగమింపజేశాం.}

ఒక విశేషమైన ఫలితం కిందిది:

\begin{thm}\ollabel{thm:universal-tm} ఇన్‌పుట్ $\tuple{e,n}$పై ప్రారంభించినప్పుడు
  \begin{enumerate}
    \item $M_e$ ఇన్‌పుట్~$n$పై ఆగితే, అప్పుడే మాత్రమే ఆగే, మరియు
    \item $M_e$ అవుట్‌పుట్ $m$తో ఆగితే, తానూ అదే అవుట్‌పుట్‌తో ఆగే
  \end{enumerate}
  ఒక \emph{సార్వత్రిక ట్యూరింగ్ యంత్రం}~$U$ ఉంది. అందువల్ల $M_e$ను
  ఇన్‌పుట్~$n$పై ప్రారంభించినప్పుడు అవుట్‌పుట్~$m$తో ఆగితే
  $f(e,n) = m$గా, లేకపోతే నిర్వచితం కాకుండా ఉండే
  $f\colon \Nat \times \Nat \pto \Nat$ ప్రమేయాన్ని $U$ గణిస్తుంది.
\end{thm}

\begin{proof}
  $U$ అత్యంత సంక్లిష్టమైన యంత్రం కనుక, దాన్ని నిజంగా నిర్మించడం దాదాపు అసాధ్యం.
  కానీ అది ఎలా పనిచేస్తుందో స్థూలంగా వర్ణించి, తరువాత చర్చ్--ట్యూరింగ్
  సిద్ధాంతప్రతిపాదనను ఆశ్రయించవచ్చు. ప్రారంభమైనప్పుడు $U$ టేపుపై
  $e$~$\TMstroke$ల దిమ్మ, దాని తరువాత $n$~$\TMstroke$ల దిమ్మ ఉంటాయి.
  అది మొదట ఇన్‌పుట్~$n$కు కుడివైపున సూచిక~$e$ను ``విసంకేతీకరిస్తుంది.''
  దీనివల్ల యంత్రం~$M_e$ యొక్క నిర్దేశాలను వర్ణించే సంఖ్యల జాబితా వస్తుంది
  (అంటే, $\TMblank$లతో వేరుచేసిన $\TMstroke$ల దిమ్మలు). తరువాత $M_e$
  వర్ణనకు కుడివైపున $M_e$ ప్రారంభ స్థితి సంఖ్యనూ, సంఖ్య~$1$నూ $U$
  వ్రాస్తుంది. (మళ్లీ, ఇవి ఏకాంక రూపంలో $\TMstroke$ల దిమ్మలుగా
  సూచించబడతాయి.) తరువాత ఇన్‌పుట్‌ను ($n$~$\TMstroke$ల దిమ్మను) కుడివైపునకు
  నకలు చేస్తుంది---కానీ ప్రతి $\TMstroke$ను మూడు $\TMstroke$ల దిమ్మతో
  మారుస్తుంది ($\TMstroke$ సంకేతపు సంఖ్య~$3$ అని గుర్తుంచుకోండి;
  $\TMendtape$ సంఖ్య~$1$, $\TMblank$ సంఖ్య~$2$). ఈ దిమ్మల క్రమం ఎడమ చివర
  ($U$ టేపుపై అవి $\TMblank$ సంకేతాలతో వేరుచేయబడ్డాయి) $\TMendtape$కు
  సంకేతమైన ఒకే~$\TMstroke$ను వ్రాస్తుంది.

  ఇప్పుడు $U$ టేపుపై ఇవి ఉన్నాయి: సూచిక~$e$, సంఖ్య~$n$, ప్రారంభ స్థితి
  సంకేత సంఖ్య (``ప్రస్తుత స్థితి''), ప్రారంభ శీర్ష స్థానం సంఖ్య~$1$
  (``ప్రస్తుత శీర్ష స్థానం''), మరియు ``టేపు'' యొక్క ప్రారంభ విషయం
  ($M_e$ సంకేతాల సంకేత సంఖ్యలకు ప్రాతినిధ్యం వహించే, $\TMblank$లతో
  వేరుచేసిన $\TMstroke$ల దిమ్మల క్రమం---అవి ``సంకేతాలు'').

  ఇన్‌పుట్~$n$పై ప్రారంభించినప్పుడు $M_e$ చేసే పనిని ఇప్పుడు $U$ కింది
  విధంగా అనుకరిస్తుంది:
  \begin{enumerate}
    \item ``ప్రస్తుత శీర్ష స్థానం'' సంఖ్య~$k$ను కనుగొనడం (మొదట అది~$1$),
    \item అక్కడి ``సంకేతం'' ఏమిటో చూడటానికి ``టేపు''లోని $k$వ దిమ్మకు వెళ్లడం,
    \item\ollabel{find-inst}%
    ప్రస్తుత ``స్థితి'', ``సంకేతం''తో సరిపోలే నిర్దేశాన్ని కనుగొనడం,
    \item ``టేపు''లోని $k$వ దిమ్మకు తిరిగి వెళ్లి, అక్కడి ``సంకేతాన్ని''
    $M_e$ వ్రాసే సంకేతపు సంకేత సంఖ్యతో మార్చడం,
    \item ప్రస్తుత ``స్థితి''ని నమోదు చేసే చోటికి శీర్షాన్ని తరలించి, అక్కడి
    సంఖ్యను కొత్త స్థితి సంఖ్యతో మార్చడం,
    \item ``టేపు స్థానం''ను నమోదు చేసే చోటికి వెళ్లి, నిర్దేశం వరుసగా ఎడమకు
    లేదా కుడికి కదలమంటే ఒక~$\TMstroke$ను చెరపడం లేదా ఒక~$\TMstroke$ను చేర్చడం.
    \item పునరావృతం చేయడం.\footnote{ఇక్కడ కొన్ని సూక్ష్మమైన ఇబ్బందులను
    దాటవేస్తున్నాం. ఉదాహరణకు, ``ప్రస్తుత శీర్ష స్థానం''ను నమోదు చేసే
    గణకాన్ని పెంచినప్పుడు $U$కు అదనపు స్థలం అవసరం కావచ్చు---ఆ సందర్భంలో
    మొత్తం ``టేపు''ను అది కుడివైపుకు తరలించాలి.}
  \end{enumerate}
  ఇన్‌పుట్~$n$పై ప్రారంభించిన $M_e$ ఎప్పటికీ ఆగకపోతే, $U$ కూడా ఎప్పటికీ
  ఆగదు; కాబట్టి దాని అవుట్‌పుట్ నిర్వచితం కాదు.

  \olref{find-inst} అంచెలో $M_e$ వర్ణనలో ప్రస్తుత ``స్థితి''/``సంకేతం''
  యుగ్మానికి నిర్దేశం లేదని తేలితే, $M_e$ ఆగేది. ఇలా జరిగితే, $U$ తన
  టేపులోని ``టేపు''కు ఎడమవైపున ఉన్న భాగాన్ని చెరిపేస్తుంది. మూడు
  $\TMstroke$ల ప్రతి దిమ్మకూ ($M_e$ టేపుపై ఒక~$\TMstroke$కు ప్రాతినిధ్యం
  వహించేది), అది తన సొంత టేపు ఎడమ చివర ఒక~$\TMstroke$ను వ్రాసి,
  ``టేపు''ను వరుసగా చెరిపేస్తుంది. ఇది పూర్తయినప్పుడు, $U$ టేపుపై
  పొడవు~$m$ గల ఒకే $\TMstroke$ల దిమ్మ ఉంటుంది.

  ``టేపు''పై మూడు~$\TMstroke$ల దిమ్మ కాని దేన్నైనా $U$ ఎదుర్కొంటే, అది
  వెంటనే ఆగుతుంది. ఈ సందర్భంలో $U$ టేపుపై ఒకే~$\TMstroke$ల దిమ్మ ఉండదు
  కనుక, దాని అవుట్‌పుట్ సహజ సంఖ్య కాదు; అంటే ఈ సందర్భంలో $f(e,n)$ నిర్వచితం
  కాదు.
\end{proof}

\end{document}
